\clearpage
\VerdierPhysicalPageAnchor{0039}%
\VerdierSemanticAnchor{ch1-1-2-8}
\noindent\textbf{定义 1.2.8.}\quad
设 $(\mathcal C_e,G^\circ,T,c)$ 与
$(\mathcal C'_e,G^\circ,T',c')$ 为两个 $G^\circ$-范畴，并设
$(F_e,m)$ 与 $(F'_e,m')\colon\mathcal C_e\to\mathcal C'_e$ 为两个
$G^\circ$-函子。若对每个 $g\in G$，图 (1.2.7.2) 都交换，就称函子态射
$u_e\colon F_e\to F'_e$ \emph{与 $G^\circ$ 的作用相容}。

\VerdierSemanticAnchor{ch1-1-2-additive}
最后，若一个 $G^\circ$-范畴的底范畴是加性范畴，就称它是\emph{加性的}。
由 (1.1.3)，与一个加性分次范畴相伴的 $G^\circ$-范畴是加性
$G^\circ$-范畴。

\VerdierSemanticAnchor{ch1-1-3}
\subsection{两种语言的等价}

\VerdierSemanticAnchor{ch1-1-3-1}
\noindent\textbf{1.3.1.}\quad
固定一个宇宙 $U$。以 $\operatorname{Grad}(G)$ 表示这样一个 $U$-范畴：其对象
是属于 $U$ 的 $G$-分次范畴，其态射是分次函子。同样，以
$G^\circ\text{-}\operatorname{Cat}$ 表示这样一个 $U$-范畴：其对象是属于
$U$ 的 $G^\circ$-范畴，其态射是 $G^\circ$-函子。由 (1.2.1)、(1.2.3) 与
(1.2.5) 立即可知，为每个 $G$-分次范畴选择一个正规劈分，便定义了一个函子

\VerdierSemanticAnchor{ch1-1-3-1-eq1}
\begin{equation}
  \Psi\colon \operatorname{Grad}(G)
  \longrightarrow G^\circ\text{-}\operatorname{Cat}.
  \tag{1.3.1.1}
\end{equation}

[18] 的结果可以解释如下。

\VerdierSemanticAnchor{ch1-1-3-2}
\noindent\textbf{命题 1.3.2.}\quad
\emph{函子 $\Psi$ 是范畴的一个等价。}

\VerdierSemanticAnchor{ch1-1-3-2-quasi-inverse}
$\Psi$ 的一个拟逆函子的描述可见 [18]。这里我们只迅速描述与一个
$G^\circ$-范畴 $(\mathcal C,G^\circ,T,c)$ 相伴的分次范畴 $\mathcal C[G]$。
首先，$\mathcal C[G]$ 的对象集就是 $\mathcal C$ 的对象集。设 $X,Y$ 为
$\mathcal C$ 的两个对象。置

\VerdierSemanticAnchor{ch1-1-3-2-eq1}
\begin{equation}
  \operatorname{Hom}^{g}(X,Y)
  =\operatorname{Hom}_{\mathcal C}(X,T(g)Y),
  \qquad g\in G,
  \tag{1.3.2.1}
\end{equation}

\VerdierSemanticAnchor{ch1-1-3-2-eq2}
\begin{equation}
  \operatorname{Hom}_{\mathcal C[G]}(X,Y)
  =\coprod_{g\in G}\operatorname{Hom}^{g}(X,Y).
  \tag{1.3.2.2}
\end{equation}

\VerdierSemanticAnchor{ch1-1-3-2-composition}
现在描述 $\mathcal C[G]$ 中态射的复合。设 $X,Y,Z$ 为 $\mathcal C$ 的三个对象，
并设 $u\in\operatorname{Hom}^{g}(X,Y)$、
$v\in\operatorname{Hom}^{g'}(Y,Z)$，其中 $g,g'\in G$。函子 $T(g)$ 定义一个映射

\VerdierSemanticAnchor{ch1-1-3-2-eq3}
\begin{equation}
  \operatorname{Hom}^{g'}(Y,Z)
  \longrightarrow
  \operatorname{Hom}_{\mathcal C}
    \bigl(T(g)Y,T(g)T(g')Z\bigr)
  \tag{1.3.2.3}
\end{equation}
